\section{Auswertung}

\subsection{Eichung des Magneten}
Die gemessenen Werte zur Eichung des Magneten sind in Tabelle \ref{tab:tabe1} angegeben.
\begin{table}[H]
  \centering
  \caption{Messwerte der Wärmepumpe}
  \label{tab:tabe1}
    \begin{tabular}{S S S S S S}
    \toprule
    $ t  \: / \si{\second} $ & $ p_a \: / \si{\bar} $ & $ p_b \: / \si{\bar} $ &
    $ T_1 \: / \si{\kelvin} $ & $ T_2 \: / \si{\kelvin} $ & $ P \: / \: \si{\watt} $\\
    \midrule
    0 & 5.0 & 5.0 & 293.65 & 293.65 & 0 \\
    60 & 4.7 & 6.0 & 294.15 & 293.55 & 115 \\
    120 & 4.4 & 6.4 & 295.15 & 293.15 & 118 \\
    180 & 4.5 & 6.9 & 296.35 & 291.95 & 122 \\
    240 & 4.6 & 7.0 & 297.55 & 290.95 & 125 \\
    300 & 4.6 & 7.0 & 298.85 & 289.95 & 125 \\
    360 & 4.5 & 7.2 & 300.05 & 289.15 & 123 \\
    420 & 4.4 & 7.4 & 301.15 & 288.45 & 123 \\
    480 & 4.3 & 7.8 & 302.35 & 287.65 & 122 \\
    540 & 4.2 & 8.0 & 303.55 & 286.95 & 122 \\
    600 & 4.2 & 8.1 & 304.65 & 286.25 & 121 \\
    660 & 4.1 & 8.3 & 305.75 & 285.55 & 121 \\
    720 & 4.0 & 8.5 & 306.75 & 284.95 & 121 \\
    780 & 4.0 & 8.8 & 307.75 & 284.35 & 121 \\
    840 & 3.9 & 9.0 & 308.75 & 283.75 & 121 \\
    900 & 3.8 & 9.1 & 309.65 & 283.15 & 121 \\
    960 & 3.8 & 9.2 & 310.55 & 282.55 & 122 \\
    1020 & 3.8 & 9.5 & 311.45 & 282.05 & 122 \\
    1080 & 3.7 & 9.8 & 312.25 & 281.55 & 122 \\
    1140 & 3.7 & 10.0 & 313.05 & 281.15 & 122 \\
    1200 & 3.7 & 10.0 & 313.9 & 280.65 & 122 \\
    1260 & 3.6 & 10.2 & 314.65 & 280.25 & 123 \\
    1320 & 3.6 & 10.3 & 315.35 & 279.85 & 123 \\
    1380 & 3.6 & 10.6 & 316.15 & 279.45 & 124 \\
    1440 & 3.6 & 10.8 & 316.85 & 279.15 & 124 \\
    1500 & 3.6 & 11.0 & 317.55 & 278.75 & 124 \\
    1560 & 3.6 & 11.1 & 318.25 & 278.55 & 124 \\
    1620 & 3.6 & 11.2 & 318.95 & 278.25 & 125 \\
    1680 & 3.5 & 11.4 & 319.55 & 277.95 & 125 \\
    1740 & 3.5 & 11.5 & 320.15 & 277.65 & 125 \\
    1800 & 3.5 & 11.7 & 320.75 & 277.45 & 125 \\
    1860 & 3.5 & 11.9 & 321.35 & 277.25 & 125 \\
    1920 & 3.5 & 12.0 & 321.95 & 277.05 & 125 \\
    1980 & 3.5 & 12.1 & 322.45 & 276.95 & 125 \\








      \bottomrule
    \end{tabular}
\end{table}

Um eine Eichkurve zu bestimmen werden diese Werte mit einem Polynom dritten Grades
der Form
\begin{equation}
  B(I)=a\cdot I^3+b\cdot I^2 + c \cdot I +d
\end{equation}
gefittet.
Die resultierenden Parameter lauten
\begin{align*}
  a &= \SI{-0.087(7)}{\milli\tesla\per\cubic\ampere} \\
  b &= \SI{1.85(18)}{\milli\tesla\per\square\ampere}\\
  c &= \SI{48.7(15)}{\milli\tesla\per\ampere}\\
  d &= \SI{10(3)}{\milli\tesla} \: .
\end{align*}
Die entsprechende Kurve ist zusammen mit den Messwerten in Abbildung \ref{fig:plot1}
dargestellt.
\begin{figure}
  \centering
  \includegraphics[height=9cm]{Plot1.pdf}
  \caption{Eichkurve und Messwerte des Magneten}
  \label{fig:plot1}
\end{figure}
Somit ergibt sich die Magnetfeldstärke bei gegebenem Strom über die Formel
\begin{equation}
  B(I)=\SI{-0.087}{\milli\tesla\per\cubic\ampere}\cdot I^3
  +\SI{1.85}{\milli\tesla\per\square\ampere}\cdot I^2
  +\SI{48.7}{\milli\tesla\per\ampere} \cdot I + \SI{10}{\milli\tesla} \; ,
  \label{eqn:bfeld}
\end{equation}
wobei der Fehler sich gemäß der Gaußschen Fehlerfortpflanzung
\begin{equation}
  \increment f = \sqrt{ \sum_{i=1}^N \left( \frac{\partial f}{\partial x_i}\right)^2
  \cdot (\increment x_i)^2  }
  \label{eqn:gaus}
\end{equation}
über
\begin{equation}
  \begin{split}
  \increment B(I) &= \sqrt{ \left( \increment a \cdot I^3 \right)^2 + \left( \increment b \cdot I^2 \right)^2
  + \left( \increment c \cdot I \right)^2 + \left( \increment d \right)^2} \\\
  &= \sqrt{ \left( \SI{0.007}{\milli\tesla\per\cubic\ampere} \cdot I^3 \right)^2 +
  \left(  \SI{0.18}{\milli\tesla\per\square\ampere} \cdot I^2 \right)^2
  + \left( \SI{1.5}{\milli\tesla\per\ampere} \cdot I \right)^2 +
  \left( \SI{3}{\milli\tesla} \right)^2}
  \label{eqn:fbfeld}
\end{split}
\end{equation}
ergibt.
\subsection{Auswertung der roten Linie}
Die rote Linie bei $\SI{643.8}{\nano\meter}$  wurde einmal ohne Magnetfeld gemessen und einmal
bei einem Magnetfeld mit einem Strom von $\SI{10}{\ampere}$. Nach Gleichungen
\ref{eqn:bfeld} und \ref{eqn:fbfeld} entspricht dies einer Magnetfeldstärke von $\SI{0.595(25)}{\tesla}$,
wobei sich der Theoriewert $g_J\,=\,1,0$ aus Gleichung \ref{eqn:gJ} ergibt und allgemein für den normalen Zeeman
Effekt gilt, um den es sich hier handelt.
Die aufgenommnen Bilder sind in Abbildungen \ref{fig:rot1} und \ref{fig:rot2} angegeben.
\begin{figure}
\begin{subfigure}[c]{0.5\textwidth}

\includegraphics[width=0.8\textwidth]{Rot1.JPG}
\subcaption{\small Linien ohne Aufspaltung}
\label{fig:rot1}
\end{subfigure}
\begin{subfigure}[c]{0.5\textwidth}
\includegraphics[width=0.8\textwidth]{Rot2.JPG}
\subcaption{\small Linien mit Aufspaltung}
\label{fig:rot2}
\end{subfigure}
\caption{Bilder der roten Cadmium Linie}
\end{figure}
Aus den Bildern werden die Abstände der Linien in Pixeln vermessen. Die daraus resultierenden Messwerte sind
in Tabelle \ref{tab:tabe2} angegeben.
\begin{table}[H]
  \centering
  \caption{Wertetabelle für $\alpha$ und $C_V$.}
  \label{tab:tab2}
    \begin{tabular}{S S S S S}
    \toprule
    $ T\: \text{in}\: \si{\K} $ & $ {\alpha \cdot 10^{-6} \: \text{in}\: \si {\per\K}} $ &
    $ C_V \: \text{in}\: \si{\J\per\K\mol} $\\
    \midrule %Cv, a *10-6, Cv
    %0 & 1 & 1\\
    88.60\pm0.24 & 9.56\pm0.06 & 14.17\pm8.13  \\ %&3.6 & 318.97\pm0.85\\
    93.81\pm0.24 & 10.10\pm0.06 & 17.58\pm10.03 \\ %& 4.7 & 440.90\pm1.11\\
    99.74\pm0.24 & 10.66\pm0.05 & 15.52\pm8.84 \\ %& 5.1 & 508.68\pm1.21\\
    104.74\pm0.24 & 11.07\pm0.05 & 18.44\pm10.52 \\ %& 4.6 & 481.79\pm1.09\\
    110.94\pm0.24 &  11.54\pm0.05 & 14.86\pm8.45 \\ %& 5.3 & 587.97\pm1.27\\
    115.96\pm0.24 & 11.89\pm0.05 & 18.49\pm10.52 \\ %& 4.6 & 533.41\pm1.10\\
    121.47\pm0.24 &  12.22\pm0.05 & 16.83\pm9.57 \\ %& 4.9 & 595.21\pm1.17\\
    126.99\pm0.24 & 12.53\pm0.04 & 16.79\pm9.54 \\ %& 4.9 & 622.29\pm1.18\\
    131.58\pm0.24 & 12.77\pm0.04 & 20.42\pm11.62 \\ %& 4.2 & 552.62\pm1.01\\
    136.65\pm0.24 & 13.02\pm0.04 & 18.40\pm10.47 \\ %& 4.6 & 628.57\pm1.11\\
    141.49\pm0.24 & 13.24\pm0.04 & 19.28\pm10.97 \\ %& 4.4 & 622.54\pm1.07\\
    146.34\pm0.24 & 13.44\pm0.04 & 19.24\pm10.95 \\ %& 4.4 & 643.88\pm1.07\\
    150.95\pm0.24 & 13.62\pm0.04 & 20.22\pm11.52 \\ %& 4.3 & 649.11\pm1.05\\
    155.34\pm0.24 & 13.79\pm0.04 & 21.31\pm12.14 \\ %& 4.1 & 636.88\pm0.98\\
    159.97\pm0.24 & 13.95\pm0.04 & 20.12\pm11.47 \\ %& 4.3 & 687.89\pm1.05\\
    164.62\pm0.24 & 14.10\pm0.04 & 20.18\pm11.51 \\ %& 4.3 & 707.87\pm1.06\\
    168.79\pm0.25 & 14.23\pm0.04 & 22.54\pm12.86 \\ %& 3.9 & 658.27\pm0.95\\
    173.45\pm0.25 &  14.37\pm0.04 & 20.08\pm11.46 \\ %& 4.3 & 745.84\pm1.06\\
    178.13\pm0.25 &  14.50\pm0.04 & 20.04\pm11.44 \\ %& 4.3 & 765.94\pm1.06\\
    182.56\pm0.25 &  14.62\pm0.04 & 21.11\pm12.06\\
    192.70\pm0.25 &  14.87\pm0.04 & 18.41\pm10.47\\
    200.15\pm0.25 &  15.04\pm0.04 & 25.19\pm14.28\\
    208.87\pm0.25 &  15.23\pm0.04 & 21.43\pm12.18\\
    217.12\pm0.25 &  15.38\pm0.04 & 22.65\pm12.88\\
    225.15\pm0.25 &  15.53\pm0.03 & 23.27\pm13.24\\
    232.70\pm0.25 &  15.70\pm0.03 & 24.75\pm14.08\\
    240.53\pm0.25 &  15.74\pm0.03 & 23.84\pm13.58\\
    248.39\pm0.25 &  15.89\pm0.03 & 23.74\pm13.53& \\
    256.01\pm0.25 &  15.97\pm0.03 & 24.46\pm13.94 \\
    263.41\pm0.26 &  16.01\pm0.03 & 25.22\pm14.38 \\
    271.08\pm0.26 &  16.18\pm0.03 & 24.26\pm13.86 \\
    278.52\pm0.26 &  16.27\pm0.03 & 25.03\pm14.29&\\
    285.98\pm0.26 &  16.35\pm0.03 & 24.92\pm14.25 \\
    293.21\pm0.26 &  16.42\pm0.03 & 25.74\pm14.72 \\
    300.98\pm0.26 &  16.50\pm0.03 & 23.87\pm13.68 \\
    308.51\pm0.26 &  16.57\pm0.03 & 24.63\pm14.12\\



      \bottomrule
    \end{tabular}
\end{table}

Die entsprechenden Wellenlängenänderungen lassen sich über
\begin{equation}
  \delta \lambda = \frac{\increment S}{2 \delta S} \increment \lambda_D
\end{equation}
berechnen, wobei $\increment \lambda_D$ für die rote Linie bei $\SI{643.8}{\nano\meter}$
einem Wert von $\SI{0.0489}{\nano\meter}$ entspricht.
\begin{table}[H]
  \centering
  \caption{Wertetabelle für $\alpha$ und $C_V$.}
  \label{tab:tab2}
    \begin{tabular}{S S S S S}
    \toprule
    $ T\: \text{in}\: \si{\K} $ & $ {\alpha \cdot 10^{-6} \: \text{in}\: \si {\per\K}} $ &
    $ C_V \: \text{in}\: \si{\J\per\K\mol} $\\
    \midrule %Cv, a *10-6, Cv
    %0 & 1 & 1\\
    88.60\pm0.24 & 9.56\pm0.06 & 14.17\pm8.13  \\ %&3.6 & 318.97\pm0.85\\
    93.81\pm0.24 & 10.10\pm0.06 & 17.58\pm10.03 \\ %& 4.7 & 440.90\pm1.11\\
    99.74\pm0.24 & 10.66\pm0.05 & 15.52\pm8.84 \\ %& 5.1 & 508.68\pm1.21\\
    104.74\pm0.24 & 11.07\pm0.05 & 18.44\pm10.52 \\ %& 4.6 & 481.79\pm1.09\\
    110.94\pm0.24 &  11.54\pm0.05 & 14.86\pm8.45 \\ %& 5.3 & 587.97\pm1.27\\
    115.96\pm0.24 & 11.89\pm0.05 & 18.49\pm10.52 \\ %& 4.6 & 533.41\pm1.10\\
    121.47\pm0.24 &  12.22\pm0.05 & 16.83\pm9.57 \\ %& 4.9 & 595.21\pm1.17\\
    126.99\pm0.24 & 12.53\pm0.04 & 16.79\pm9.54 \\ %& 4.9 & 622.29\pm1.18\\
    131.58\pm0.24 & 12.77\pm0.04 & 20.42\pm11.62 \\ %& 4.2 & 552.62\pm1.01\\
    136.65\pm0.24 & 13.02\pm0.04 & 18.40\pm10.47 \\ %& 4.6 & 628.57\pm1.11\\
    141.49\pm0.24 & 13.24\pm0.04 & 19.28\pm10.97 \\ %& 4.4 & 622.54\pm1.07\\
    146.34\pm0.24 & 13.44\pm0.04 & 19.24\pm10.95 \\ %& 4.4 & 643.88\pm1.07\\
    150.95\pm0.24 & 13.62\pm0.04 & 20.22\pm11.52 \\ %& 4.3 & 649.11\pm1.05\\
    155.34\pm0.24 & 13.79\pm0.04 & 21.31\pm12.14 \\ %& 4.1 & 636.88\pm0.98\\
    159.97\pm0.24 & 13.95\pm0.04 & 20.12\pm11.47 \\ %& 4.3 & 687.89\pm1.05\\
    164.62\pm0.24 & 14.10\pm0.04 & 20.18\pm11.51 \\ %& 4.3 & 707.87\pm1.06\\
    168.79\pm0.25 & 14.23\pm0.04 & 22.54\pm12.86 \\ %& 3.9 & 658.27\pm0.95\\
    173.45\pm0.25 &  14.37\pm0.04 & 20.08\pm11.46 \\ %& 4.3 & 745.84\pm1.06\\
    178.13\pm0.25 &  14.50\pm0.04 & 20.04\pm11.44 \\ %& 4.3 & 765.94\pm1.06\\
    182.56\pm0.25 &  14.62\pm0.04 & 21.11\pm12.06\\
    192.70\pm0.25 &  14.87\pm0.04 & 18.41\pm10.47\\
    200.15\pm0.25 &  15.04\pm0.04 & 25.19\pm14.28\\
    208.87\pm0.25 &  15.23\pm0.04 & 21.43\pm12.18\\
    217.12\pm0.25 &  15.38\pm0.04 & 22.65\pm12.88\\
    225.15\pm0.25 &  15.53\pm0.03 & 23.27\pm13.24\\
    232.70\pm0.25 &  15.70\pm0.03 & 24.75\pm14.08\\
    240.53\pm0.25 &  15.74\pm0.03 & 23.84\pm13.58\\
    248.39\pm0.25 &  15.89\pm0.03 & 23.74\pm13.53& \\
    256.01\pm0.25 &  15.97\pm0.03 & 24.46\pm13.94 \\
    263.41\pm0.26 &  16.01\pm0.03 & 25.22\pm14.38 \\
    271.08\pm0.26 &  16.18\pm0.03 & 24.26\pm13.86 \\
    278.52\pm0.26 &  16.27\pm0.03 & 25.03\pm14.29&\\
    285.98\pm0.26 &  16.35\pm0.03 & 24.92\pm14.25 \\
    293.21\pm0.26 &  16.42\pm0.03 & 25.74\pm14.72 \\
    300.98\pm0.26 &  16.50\pm0.03 & 23.87\pm13.68 \\
    308.51\pm0.26 &  16.57\pm0.03 & 24.63\pm14.12\\



      \bottomrule
    \end{tabular}
\end{table}

Durch die Gleichung
\begin{equation}
  \bar{x} = \frac{1}{N} \sum_{i=1}^{N} x_i \: \:
  \label{eqn:mit}
\end{equation}
\noindent lässt sich der Mittelwert bilden, wobei der dazugehörige Fehler sich durch
\begin{equation}
  \increment \bar{x} = \frac{1}{\sqrt{N}} \sqrt{ \frac{1}{N-1} \sum_{i=1}^N
  (x_i - \bar{x})^2}
  \label{eqn:mitf}
\end{equation}
ergibt. Somit ergibt sich für $\delta \lambda$ ein Mittelwert von
$ \delta \lambda = \SI{11.52(20)}{\pico\meter} $.
Um den Faktor $g_J$ zu bestimmen, muss zunächst die Energiedifferenz $\increment E$ bestimmt werden, welche
der Wellenlängenänderung $\increment \lambda$ entspricht. Zu diesem Zweck wird die Formel
\begin{equation}
  E=\frac{\text{h}\, \text{c} }{\lambda}
\end{equation}
mit dem Plank´schen Wirkungsquantum h und der Lichtgeschwindigkeit c um kleine Änderungen $\increment \lambda$
in erster Ordnung entwickelt, sodass sich
\begin{equation}
  \increment E =\frac{\text{h}\, \text{c} }{\lambda^2}\increment \lambda
  \label{eqn:de}
\end{equation}
ergibt
mit dem Fehler
\begin{equation}
  \delta E =\frac{\text{h}\, \text{c} }{\lambda^2}\delta \lambda
  \label{eqn:fde}
\end{equation}
durch Gleichung \ref{eqn:gaus}.
Zusammen mit Gleichung \ref{eqn:emag} ergibt sich letztendlich
\begin{align*}
g_J=1,00 \pm 0,04 \: .
\end{align*}

\subsection{Auswertung der blauen Linie}
Für die blaue Linie bei $\SI{480}{\nano\meter}$  wurde neben der Aufnahme ohne Magnetfeld einmal
die Aufspaltung bei einem Strom von $\SI{5.5}{\ampere}$ für die $\sigma$-Linien aufgenommen
und einmal bei einem Strom von $\SI{14}{\ampere}$ für die $\pi$-Linien. Nach Gleichungen
\ref{eqn:bfeld} und \ref{eqn:fbfeld} entspricht dies Magnetfeldstärken von $\SI{0.319(10)}{\tesla}$
und $\SI{0.82(5)}{\tesla}$.
Die aufgenommnen Bilder sind in Abbildungen \ref{fig:blau1}, \ref{fig:blau2} und \ref{fig:blau3} angegeben.
\begin{figure}
\begin{subfigure}[c]{0.5\textwidth}

\includegraphics[width=0.8\textwidth]{Blau1.JPG}
\subcaption{\small Linien ohne Aufspaltung}
\label{fig:blau1}
\end{subfigure}
\begin{subfigure}[c]{0.5\textwidth}
\includegraphics[width=0.8\textwidth]{Blau2.JPG}
\subcaption{\small $\sigma$-Linien mit Aufspaltung}
\label{fig:blau2}
\end{subfigure}
\begin{subfigure}[c]{0.5\textwidth}
\includegraphics[width=0.8\textwidth]{Blau3.JPG}
\subcaption{\small $\pi$-Linien mit Aufspaltung}
\label{fig:blau3}
\end{subfigure}
\caption{Bilder der blauen Cadmium Linie}
\end{figure}
Die gemessenen Abstände und die entsprechenden Wellenlängenänderungen mit dem Wert
$\increment \lambda_D = \SI{0.02695}{\nano\meter}$ der blauen Linie sind in Tabelle
\ref{tab:tabe3} angegeben.
\begin{table}
  \centering
  \caption{Messwerte für den ersten Doppelspalt.}
   \begin{tabular}{S S| S S | S S}
    \toprule
    $x/\; \si{\mm}$& $A/\;\si{\nA}$ &
    $x/\; \si{\mm}$& $A/\;\si{\nA}$ &
    $x/\; \si{\mm}$& $A/\;\si{\nA}$ \\
    \midrule

    15.0& 4.6& 23.0& 25.0& 29.5& 6.0\\
    15.5& 4.2& 23.5& 30.0& 30.0& 5.3\\
    16.0& 4.0& 24.0& 35.0& 30.5& 4.9\\
    16.5& 4.0& 24.25& 36.0& 31.0& 4.7\\
    17.0& 4.4& 24.5& 37.0& 31.5& 4.4\\
    17.5& 5.5& 24.75& 38.0& 32.0& 4.2\\
    18.0& 6.6& 25.00& 37.0& 32.5& 3.8\\
    18.5& 7.7& 25.25& 36.0& 33.0& 3.6\\
    19.0& 8.2& 25.5& 36.0& 33.5& 3.2\\
    19.5& 8.4& 26.0& 33.0& 34.0& 3.2\\
    20.0& 8.4& 26.5& 28.5& 34.5& 3.2\\
    20.25& 8.4& 27.0& 23.0& 35.0& 3.3\\
    20.5& 8,7& 27.5& 18.0& 35.5& 3.4\\
    21.0& 9.8& 28.0& 13.5& 36.0& 3.5\\
    21.5& 12.0& 28.5& 10.0\\
    22.0& 15.0& 29.0& 7.8\\
    22.5& 20.0& 29.25& 6.7\\


   \bottomrule
  \end{tabular}
  \label{tab:tabelle3}
\end{table}

Der Mittelwert der Wellenlängenänderung ergibt sich nach Gleichungen \ref{eqn:mit} und
\ref{eqn:mitf} zu $ \delta \lambda_{\sigma} = \SI{6.17(4)}{\pico\meter} $ für die $\sigma$-Linie
und zu $ \delta \lambda_{\pi} = \SI{5.01(5)}{\pico\meter} $ für die $\pi$-Linie.

Durch die Gleichungen \ref{eqn:de}, \ref{eqn:fde} und \ref{eqn:emag} ergibt sich somit
\begin{align}
  g_J&=1,80 \pm 0,06 \; , \sigma-\text{Linie}\\
  g_J&=0,57 \pm 0,032 \; , \pi-\text{Linie}
\end{align}
für die Werte von $g_J$ der beiden Linien.
